\ac{ERP}-Systeme unterstützen die unternehmensweite Planung, Steuerung und Auswertung von Geschäftsprozessen und Ressourcen und bilden damit eine zentrale Komponente moderner Unternehmens-IT \cite[\pagef 6]{nickel_wie_2020}.

Vor dem Hintergrund der fortschreitenden Digitalisierung sowie der kontinuierlichen Weiterentwicklung der IT unterliegen \ac{ERP}-Systeme einem stetigen Wandel und werden fortlaufend optimiert \cite[\pagef 1]{hrischev_artificial_2023}. Ein besonders bedeutender Trend hierbei ist der zunehmende Einsatz verschiedener Softwarefunktionen und -module, die auf \ac{KI} setzen. Dieser ermöglicht es, Prozesse innerhalb von \ac{ERP}-Systemen zu automatisieren und eröffnet Unternehmen neue Nutzungsmöglichkeiten. \cite[\pagef 6ff]{hrischev_artificial_2023}

Im Zentrum dieser Entwicklung stehen sogenannte \ac{KI}-Agenten, auf deren Einsatz sich diese Arbeit konzentriert. Diese übernehmen eigenständig Aufgaben, nehmen ihre Umgebung wahr, handeln auf Basis vorgegebener Rollen und passen sich durch Lernen aus Erfahrungen an. \cite[\pagef 2]{deng_ai_2025} Entsprechend können sie innerhalb eines \ac{ERP}-Systems als Assistenten fungieren, die auf Anfrage Auskünfte geben, Routineaufgaben selbstständig ausführen oder vorausschauende Analysen erstellen. Durch solche Szenarien könnte der Einsatz von \ac{KI}-Agenten eine Steigerung der Produktivität und Datenqualität sowie eine Entlastung von Routineaufgaben ermöglichen. \cite[\pagef 1033]{gujar_role_2025}